\chapter{FUNKTOR $\mathbf{\emph{K}_0}$}

Sekarang kita meninjau secara singkat apa yang disebut teori-$K$ untuk aljabar-$C^*$. Barangkali pengantar terbaik untuk pokok bahasan ini adalah~\cite{RordamLL:2000},
yang beberapa bagian elementernya diikuti dengan cukup saksama dalam catatan ini. Pengantar lain yang sangat mudah dibaca adalah~\cite{Wegge-Olsen:1993}.

Diberikan suatu aljabar-$C^*$ $A$, kita terutama akan memperhatikan dua grup yang dikenal sebagai $K_0(A)$ dan~$K_1(A)$. Bab ini membahas
grup yang pertama, yakni grup proyeksi di $A$ dengan operasi penjumlahan. Tentu saja, jika kita ingin menjumlahkan proyeksi sembarang, kita langsung
menghadapi kesulitan serius. Sebelumnya telah kita tunjukkan (dalam proposisi~\ref{prop_sum_projs}) bahwa agar proyeksi dapat dijumlahkan, proyeksi-proyeksi itu harus komutatif
(dan hasil kalinya nol)! Penyelesaian dilema ini mencerminkan cara berpikir yang lazim dalam teori-$K$ secara umum. Jika dua proyeksi tidak
komutatif, singkirkan penghalang yang mencegah keduanya menjadi komutatif. Jika tidak cukup ruang bagi keduanya untuk saling melewati, berilah keduanya lebih banyak
ruang. Jangan bersikeras memandang keduanya sebagai makhluk yang mencoba hidup dalam dunia matriks $1 \times 1$ berentri di~$A$ yang sesak tanpa harapan.
Biarkan keduanya, misalnya, menjelajahi dunia matriks $2 \times 2$ yang jauh lebih lapang. % SOURCE-CORRECTION: CH17-C001
Untuk mewujudkan gagasan ini secara teknis, kita mencari relasi ekuivalensi $\sim$ yang mengidentifikasi proyeksi $p$ di $A$ dengan matriks $\begin{bmatrix} p & \vc 0 \\ \vc 0 & \vc 0 \end{bmatrix}$
dan $\begin{bmatrix} \vc 0 & \vc 0 \\ \vc 0 & p \end{bmatrix}$. Sebagai motivasi heuristik saja---bukan sebagai konsekuensi suatu kongruensi perkalian yang telah dibuktikan---hambatan itu tampak hilang: jika $p$ dan $q$ proyeksi di~$A$, maka % SOURCE-CORRECTION: CH17-C002
    \[ pq \sim \begin{bmatrix} p & \vc 0 \\ \vc 0 & \vc 0 \end{bmatrix} \begin{bmatrix} \vc 0 & \vc 0 \\ \vc 0 & q \end{bmatrix}  \\
                   = \begin{bmatrix} \vc 0 & \vc 0 \\ \vc 0 & \vc 0 \end{bmatrix}           \\
                   =  \begin{bmatrix} \vc 0 & \vc 0 \\ \vc 0 &  q \end{bmatrix}\begin{bmatrix} p & \vc 0 \\ \vc 0 & \vc 0 \end{bmatrix}  \\
                   \sim qp\,. \]
Dalam gambaran heuristik ini, $p$ dan $q$ komutatif modulo relasi ekuivalensi $\sim$, dan kita dapat mendefinisikan penjumlahan dengan sesuatu seperti
$p \oplus q = \begin{bmatrix} p & \vc 0 \\ \vc 0 & q \end{bmatrix}$.

Cukup jelas bahwa perangkat sederhana di atas tidak akan langsung menghasilkan sesuatu yang menyerupai grup Abelian, tetapi ini merupakan suatu permulaan. Jika Anda menduga
bahwa pada akhirnya konstruksi $K_0(A)$ mungkin agak rumit, dugaan Anda sepenuhnya benar.




\section{Relasi Ekuivalensi pada Proyeksi}

Kita mendefinisikan beberapa relasi ekuivalensi, yang semuanya sesuai untuk proyeksi dalam aljabar-$C^*$ beridentitas.

\begin{defn} Dalam suatu aljabar beridentitas, elemen $a$ dan $b$ disebut
 \index{serupa!elemen suatu aljabar}%
\df{serupa} jika terdapat elemen invertibel $s$ sedemikian sehingga $b = sas^{-1}$.
 \index{<binrelequivs@$a \sim_s b$ (keserupaan elemen suatu aljabar)}%
Dalam hal ini kita menulis $a \sim_s b$.
\end{defn}

\begin{prop} Relasi keserupaan $\sim_s$ yang didefinisikan di atas merupakan relasi ekuivalensi pada elemen-elemen suatu aljabar beridentitas.
\end{prop}

\begin{defn} Dalam suatu aljabar-$C^*$, elemen $a$ dan $b$ disebut
 \index{uniter!ekuivalensi}%
\df{ekuivalen secara uniter} jika terdapat suatu elemen uniter $u \in \wt A$ sedemikian sehingga $b = uau^*$. % SOURCE-CORRECTION: CH17-C003
 \index{<binrelequivu@$a \sim_u b$ (ekuivalensi uniter)}%
Dalam hal ini kita menulis $a \sim_u b$.
\end{defn}

\begin{prop} Relasi ekuivalensi uniter $\sim_u$ yang didefinisikan di atas merupakan relasi ekuivalensi pada elemen-elemen suatu aljabar-$C^*$.
\end{prop}

\begin{prop}\label{0060224fa} Elemen $a$ dan $b$ dari suatu aljabar-$C^*$ beridentitas $A$ ekuivalen secara uniter jika dan hanya jika
terdapat elemen $u \in \ofml U(A)$ sedemikian sehingga $b = uau^*$.
\end{prop}

\begin{proof}[\emph{Petunjuk untuk bukti}] Andaikan terdapat $v \in \ofml U(\wt A)$ sedemikian sehingga $b = vav^*$. Berdasarkan proposisi~\ref{0015213},
terdapat $u \in A$ dan $\alpha \in \C$ sedemikian sehingga $v = u + \alpha \vc j$ (dengan $\vc j = \vc 1_{\wt A} - \vc 1_A$). Tunjukkan
bahwa $\abs\alpha = 1$, $u$ uniter, dan $b = uau^*$.

Untuk arah sebaliknya, andaikan terdapat $u \in \ofml U(A)$ sedemikian sehingga $b = uau^*$. Misalkan $v = u + \vc j$. \ns
\end{proof}

\begin{prop}[Dekomposisi polar]\label{0060141} Untuk setiap elemen invertibel $s$ dari suatu
 \index{dekomposisi polar!elemen invertibel aljabar-$C^*$}%
 \index{dekomposisi!polar}%
aljabar-$C^*$ beridentitas, terdapat suatu elemen uniter $\omega(s)$ dalam aljabar tersebut sedemikian sehingga % SOURCE-CORRECTION: CH17-C004
   \[   s = \omega(s) \abs s\,. \]
\end{prop}

\begin{proof}[\emph{Petunjuk untuk bukti}] Gunakan Akibat~\ref{00183313}.   \ns
\end{proof}

\begin{prop}\label{K002021} Jika $A$ suatu aljabar-$C^*$ beridentitas, maka fungsi $\omega\colon \inv A \sto \ofml U(A)$
yang didefinisikan dalam proposisi sebelumnya kontinu.
\end{prop}

\begin{proof}[\emph{Petunjuk untuk bukti}] Gunakan proposisi~\ref{000646fa} dan~\ref{001449}.   \ns
\end{proof}

\begin{prop}\label{K002024} Misalkan $s = u\abs s$ dekomposisi polar dari elemen invertibel $s$ dalam aljabar-$C^*$ beridentitas~$A$.
Maka $u \sim_h s$ dalam $\inv A$.
\end{prop}

\begin{proof}[\emph{Petunjuk untuk bukti}] Untuk $0 \le t \le 1$, misalkan $c_t = u(t\abs s + (1 - t)\vc 1)$. Simpulkan dari
proposisi~\ref{0018301} bahwa terdapat $\epsilon \in (0,1]$ sedemikian sehingga $\abs s \ge \epsilon \vc 1$. Gunakan proposisi yang sama
untuk menunjukkan bahwa $t\abs s + (1 - t)\vc 1$ invertibel untuk setiap $t \in [0,1]$.   \ns
\end{proof}

\begin{prop}\label{K002027} Misalkan $u$ dan $v$ elemen-elemen uniter dalam aljabar-$C^*$ beridentitas~$A$. Jika $u \sim_h v$ dalam $\inv A$,
maka $u \sim_h v$ dalam~$\ofml U(A)$.
\end{prop}

\begin{prop}\label{K002031} Misalkan $u_1$, $u_2$, $u_3$, dan $u_4$ elemen-elemen uniter dalam aljabar-$C^*$ beridentitas~$A$.
Jika $u_1 \sim_h u_2$ dalam $\ofml U(A)$ dan $u_3 \sim_h u_4$ dalam $\ofml U(A)$, maka $u_1u_3 \sim_h u_2u_4$ dalam~$\ofml U(A)$.
\end{prop}

\begin{exam}\label{0060116} Misalkan $h$ suatu elemen swaadjoin dari aljabar-$C^*$ beridentitas~$A$. Maka $\exp(ih)$ uniter dan
homotop dengan~$\vc 1$ dalam~$\ofml U(A)$.
\end{exam}

\begin{proof}[\emph{Petunjuk untuk bukti}] Pertimbangkan lintasan $c\colon [0,1] \sto \ofml U(A)\colon t \mapsto \exp(ith)$. Gunakan % SOURCE-CORRECTION: CH17-C005
proposisi~\ref{001449}.  \ns
\end{proof}

\begin{notn} Jika $p$ dan $q$ proyeksi dalam suatu aljabar-$C^*$ $A$, kita menulis $p \sim q$
 \index{ekuivalensi Murray--von Neumann}%
 \index{ekuivalensi!Murray--von Neumann}%
 \index{<binrelequiv@$p \sim q$ (ekuivalensi Murray--von Neumann)}%
($p$ \df{ekuivalen Murray--von Neumann} dengan~$q$) jika terdapat elemen $v \in A$
sedemikian sehingga $v^*v = p$ dan $vv^* = q$. Perhatikan bahwa $v$ demikian secara otomatis merupakan isometri
parsial. Kita akan menyebutnya
 \index{realisasi ekuivalensi Murray--von Neumann}%
 \index{ekuivalensi!Murray--von Neumann!realisasi}%
 \index{Murray--von Neumann!ekuivalensi!realisasi}%
\emph{isometri parsial yang merealisasikan ekuivalensi}.
\end{notn}

\begin{prop}\label{0060216} Relasi ekuivalensi Murray--von Neumann $\sim$ merupakan
relasi ekuivalensi pada keluarga $\fml P(A)$ dari proyeksi-proyeksi dalam suatu aljabar-$C^*$~$A$.
\end{prop}

\begin{prop}\label{0060144} Misalkan $a$ dan $b$ elemen-elemen swaadjoin dari suatu aljabar-$C^*$ beridentitas. Jika $a \sim_s b$, maka $a \sim_u b$.
Bahkan, jika $b = sas^{-1}$ dan $s = u\abs s$ adalah dekomposisi polar dari $s$, maka $b = uau^*$.
\end{prop}

\begin{proof}[\emph{Petunjuk untuk bukti}] Andaikan terdapat elemen invertibel $s$ sedemikian sehingga $b = sas^{-1}$. Misalkan $s = u\abs s$
dekomposisi polar dari~$s$. Tunjukkan bahwa $a$ komutatif dengan ${\abs s}^2$, sehingga juga komutatif dengan setiap elemen dalam aljabar $C^*(\vc 1,{\abs s}^2)$.
Secara khusus, $a$ komutatif dengan~${\abs s}^{-1}$. Dari sini diperoleh $uau^* = b$.   \ns
\end{proof}

\begin{prop}\label{0060221} Jika $p$ dan $q$ proyeksi dalam suatu aljabar-$C^*$, maka
   \[ p \sim_h q \implies p \sim_u q \implies p \sim q\,. \]
\end{prop}

\begin{proof}[\emph{Petunjuk untuk bukti}] Dalam petunjuk ini $\vc 1 = \vc 1_{\wt A}$. Untuk implikasi pertama, tunjukkan bahwa tanpa mengurangi keumuman
kita dapat mengandaikan $\norm{p - q} < \frac12$. Misalkan $s = pq + (\vc 1 - p)(\vc 1 - q)$. Buktikan bahwa $p \sim_s q$. Untuk
itu, tuliskan $s - \vc 1$ sebagai $p(q - p) + (\vc 1 - p)((\vc 1 - q) - (\vc 1 - p))$ dan gunakan akibat~\ref{cor_Neumann_series}.
Kemudian gunakan proposisi~\ref{0060144}.

Untuk membuktikan implikasi kedua, perhatikan bahwa jika $upu^* = q$ untuk suatu elemen uniter dalam~$\wt A$, maka $up$ merupakan isometri parsial
dalam~$A$.  \ns
\end{proof}

Secara umum, kebalikan dari implikasi kedua, $p \sim q \implies p \sim_u q$, dalam proposisi sebelumnya tidak berlaku (lihat % SOURCE-CORRECTION: CH17-C024
contoh~\ref{0060233} di bawah). Namun, untuk proyeksi $p$ dan $q$ dalam suatu aljabar-$C^*$ beridentitas, jika berlaku baik $p \sim q$ maupun
$\vc 1 - p \sim \vc 1 - q$, maka kita dapat menyimpulkan $p \sim_u q$.

\begin{prop}\label{0060224} Misalkan $p$ dan $q$ proyeksi dalam suatu aljabar-$C^*$ beridentitas~$A$. Maka $p \sim_u q$ jika dan hanya jika
$p \sim q$ dan $\vc 1 - p \sim \vc 1 - q$.
\end{prop}

\begin{proof}[\emph{Petunjuk untuk bukti}] Untuk arah sebaliknya, andaikan isometri parsial $v$ merealisasikan ekuivalensi $p \sim q$
dan $w$ merealisasikan $\vc 1 - p \sim \vc 1 - q$. Pertimbangkan elemen $u = v + w + \vc j$ dalam~$\wt A$.   \ns
\end{proof}

\begin{defn} Suatu elemen $s$ dari aljabar-$C^*$ beridentitas disebut
 \index{isometri!dalam aljabar-$C^*$ beridentitas}%
\df{isometri} jika $s^*s = \vc 1$.
\end{defn}

\begin{exer} Jelaskan mengapa terminologi dalam definisi sebelumnya masuk akal.
\end{exer}

Tidak sulit untuk melihat bahwa kebalikan dari implikasi kedua dalam proposisi~\ref{0060221} pada umumnya gagal.

\begin{exam}\label{0060233} Jika $p$ dan $q$ proyeksi dalam suatu aljabar-$C^*$, maka
   \[ p \sim q \,\,\,\,\not\!\!\!\!\implies p \sim_u q\,. \]
Sebagai contoh, jika $s$ merupakan isometri tak uniter (seperti geser unilateral), maka $s^*s \sim ss^*$, tetapi $s^*s \nsim_u ss^*$.
\end{exam}

\begin{proof}[\emph{Petunjuk untuk bukti}] Buktikan, dan ingatlah, bahwa tidak ada proyeksi tak nol yang dapat ekuivalen Murray--von Neumann
dengan proyeksi nol.   \ns
\end{proof}

\begin{rem} Kebalikan dari implikasi pertama, $p \sim_u q \implies p \sim_h q$, dalam proposisi~\ref{0060221} juga tidak berlaku % SOURCE-CORRECTION: CH17-C025
secara umum untuk proyeksi dalam suatu aljabar-$C^*$. Namun, contoh yang menggambarkan gejala ini tidak mudah ditemukan. Untuk melihat
apa yang terlibat, lihat~\cite{RordamLL:2000}, contoh 2.2.9 dan 11.3.4.
\end{rem}

Alangkah baiknya jika implikasi-implikasi dalam proposisi~\ref{0060221} dapat dibalik. Namun, seperti telah kita lihat, hal itu tidak terjadi.
Salah satu cara menghadapi fakta yang membandel, sebagaimana dicatat pada awal bab ini, adalah memberi objek-objek matematika
yang sedang kita kaji ruang yang lebih luas untuk bergerak. Beralihlah ke matriks. Proposisi~\ref{0060236} dan~\ref{0060244} merupakan contoh
cara kerja teknik ini. Dalam arti tertentu, kita dapat membuat implikasi-implikasi dalam~\ref{0060221} berbalik arah.

\begin{notn} Jika $a_1$, $a_2$, \dots $a_n$ elemen-elemen dari suatu aljabar-$C^*$ $A$, maka
 \index{diagonal@$\diag(a_1,\dots,a_n)$ (matriks diagonal dalam $\M nA$)}%
$\diag(a_1,a_2,\dots,a_n)$ adalah matriks diagonal dalam $\M nA$ yang diagonal utamanya terdiri atas
elemen-elemen $a_1$, \dots, $a_n$. Notasi ini juga kita gunakan untuk matriks blok. Sebagai
contoh, jika $a$ matriks $m \times m$ dan $b$ matriks $n \times n$, maka
$\diag(a,b)$ adalah matriks $(m+n) \times (m+n)$ $\begin{bmatrix} a & \vc 0 \\ \vc 0 & b
\end{bmatrix}$. (Tentu saja, $\vc 0$ di sudut kanan atas matriks ini adalah matriks $m \times n$ yang semua
entrinya nol, sedangkan $\vc 0$ di sudut kiri bawah adalah matriks $n \times m$ yang setiap entrinya nol.)
\end{notn}

\begin{prop}\label{0060236} Misalkan $p$ dan $q$ proyeksi dalam suatu aljabar-$C^*$~$A$. Maka
   \[ p \sim q \implies \diag(p,\vc 0) \sim_u \diag(q,\vc 0) \text{ dalam $\M 2A$.} \]
\end{prop}

\begin{proof}[\emph{Petunjuk untuk bukti}] Misalkan $v$ isometri parsial dalam $A$ yang merealisasikan ekuivalensi Murray--von Neumann $p \sim q$.
Pertimbangkan matriks $u = \begin{bmatrix} v  &  \vc 1_{\wt A} - q  \\  \vc 1_{\wt A} - p  &  v^* \end{bmatrix}$.  \ns
\end{proof}

Ingat dari akibat~\ref{001311005fa} bahwa spektrum setiap elemen uniter dari suatu aljabar-$C^*$ beridentitas terletak pada lingkaran satuan.
Hebatnya, jika spektrumnya bukan seluruh lingkaran, elemen itu homotop dengan~$\vc 1$ dalam~$\ofml U(A)$.

\begin{prop}\label{0060121} Misalkan $u$ suatu elemen uniter dalam aljabar-$C^*$ beridentitas. Jika $\sigma(u) \ne \T$, maka $u \sim_h \vc 1$ dalam~$\ofml U(A)$.
\end{prop}

\begin{proof}[\emph{Petunjuk untuk bukti}] Pilih $\theta \in \R$ sedemikian sehingga $\exp(i\theta) \notin \sigma(u)$. Maka terdapat fungsi kontinu yang tunggal
   \[ \phi\colon \sigma(u) \sto (\theta, \theta + 2\pi)\colon \exp(it) \mapsto t\,. \]
Misalkan $h = \phi(u)$ dan gunakan contoh~\ref{0060116}.   \ns
\end{proof}

\begin{exam}\label{0060129a} Jika $A$ suatu aljabar-$C^*$ beridentitas, maka $\begin{bmatrix} \vc 0  &  \vc 1  \\  \vc 1  &  \vc 0  \end{bmatrix} \sim_h
\begin{bmatrix} \vc 1  &  \vc 0  \\ \vc 0  & \vc 1 \end{bmatrix}$ dalam~$\ofml U(\M 2A)$.
\end{exam}

Matriks $J = \begin{bmatrix} \vc 0  &  \vc 1  \\ \vc 1 & \vc 0 \end{bmatrix}$ sangat berguna jika dipakai bersama
proposisi~\ref{K002031} untuk membentuk homotopi antara matriks-matriks dalam $\ofml U(A)$. Dalam memeriksa beberapa contoh
berikutnya, matriks ini sangat membantu. Perhatikan bahwa mengalikan matriks $2 \times 2$ dari kanan dengan $J$ menukar kolom-kolomnya; mengalikan
dari kiri dengan $J$ menukar baris-barisnya; dan mengalikan dari kiri serta kanan dengan $J$ menukar elemen-elemen pada kedua diagonal.

\begin{exam}\label{0060129b} Jika $u$ dan $v$ elemen-elemen uniter dalam suatu aljabar-$C^*$ beridentitas~$A$, maka
   \[ \diag(u,v) \sim_h \diag(v,u) \]
dalam~$\ofml U(\M 2A)$.
\end{exam}

\begin{exam}\label{0060129c} Jika $u$ dan $v$ elemen-elemen uniter dalam suatu aljabar-$C^*$ beridentitas $A$, maka
  \[ \diag(u,v) \sim_h \diag(uv, \vc 1) \]
dalam $\ofml U(\M 2A)$.
\end{exam}

\begin{exam}\label{0060129d} Jika $u$ dan $v$ elemen-elemen uniter dalam suatu aljabar-$C^*$ beridentitas~$A$, maka
   \[ \diag(uv,\vc 1) \sim_h \diag(vu, \vc 1) \]
dalam~$\ofml U(\M 2A)$.
\end{exam}

\begin{prop}\label{0060244} Misalkan $p$ dan $q$ proyeksi dalam suatu aljabar-$C^*$~$A$. Maka
   \[ p \sim_u q \implies \diag(p,\vc 0) \sim_h \diag(q,\vc 0) \text{ dalam $\ofml P(\M 2A)$.} \]
\end{prop}

\begin{exam}\label{0060247fa} Dua proyeksi dalam $\mathbf M_n$ ekuivalen Murray--von Neumann jika dan hanya jika keduanya memiliki teras yang sama,
yang pada gilirannya ekuivalen dengan jangkauan keduanya berdimensi sama.
\end{exam}

\begin{proof}[\emph{Petunjuk untuk bukti}] Gunakan proposisi~\ref{tr0121}--\ref{tr0124} dan~\ref{pi0117}--\ref{pi0121}.   \ns
\end{proof}

Dalam contoh~\ref{000533a} kita memperkenalkan aljabar matriks $n \times n$ $\M nA$, dengan $A$ suatu aljabar. Jika $A$ suatu
aljabar-$*\,$ dan $n \in \N$, kita dapat memperlengkapi $\M nA$ dengan involusi secara alami. Involusi itu didefinisikan, seperti yang dapat diduga,
sebagai analog dari ``transposisi konjugat'': jika $\vc a \in \M nA$, maka
   \[ \vc a^* = \bigl[a_{ij}\bigr]^* := \bigl[{a_{ji}}^*\bigr]. \]

Jika $A$ suatu aljabar-$C^*$, kita dapat memperkenalkan norma pada $\M nA$ yang membuatnya juga menjadi aljabar-$C^*$.
 \index{Ma@$\M nA$!sebagai aljabar-$C^*$}%
 \index{C*@aljabar-$C^*$!$\M nA$ sebagai suatu}%
Pilih representasi setia $\phi\colon A \sto \ofml B(H)$ dari $A$, dengan $H$ suatu ruang Hilbert. Untuk setiap $n \in \N$
definisikan
   \[ \phi_n\colon \M nA \sto \ofml B(H^n)\colon \bigl[a_{ij}\bigr] \mapsto \bigl[ \phi(a_{ij})\bigr] \]
dengan $H^n$ jumlah langsung $n$ rangkap dari $H$ dengan dirinya sendiri. Kemudian definisikan norma suatu elemen $\vc a \in \M nA$ sebagai
norma operator $\phi_n(\vc a) \in \ofml B(H^n)$. Yakni, $\norm{\vc a} := \norm{\phi_n(\vc a)}$.

Norma ini merupakan norma-$C^*$ pada $\M nA$. Norma tersebut tidak bergantung pada representasi khusus yang kita pilih berkat
ketunggalan norma-$C^*$ (lihat akibat~\ref{00152144}).

\begin{prop} Misalkan $A$ suatu aljabar-$C^*$ dan $\vc a \in \M nA$. Maka
   \[ \max\{\norm{a_{ij}}\colon 1 \le i,j \le n\} \le \norm{\vc a} \le \sum_{i,j = 1}^n \norm{a_{ij}}. \]
\end{prop}

\begin{proof}[\emph{Petunjuk untuk bukti}] Karena setiap representasi setia merupakan isometri (lihat proposisi~\ref{00152181}),
Anda dapat menyederhanakan persoalan dengan memandang representasi $\phi\colon A \sto \ofml B(H)$ yang dibahas di atas sebagai
pemetaan inklusi.

Untuk ketaksamaan pertama, bagi $v \in H$ definisikan vektor $\vc E_j(v) \in H^n$ sebagai $n$-tupel
yang semua entrinya nol kecuali entri ke-$j$, yang bernilai~$v$. Kemudian periksa bahwa
   \[ \norm{a_{ij}v} \le \norm{\vc a\vc E_j(v)} \le \norm{\vc a} \]
setiap kali $\vc a \in \M nA$, $v \in H$, dan $\norm v \le 1$.

Untuk ketaksamaan kedua, tunjukkan bahwa
   \[ {\norm{\vc a\vc v}}^2 \le \sum_{i=1}^n\biggl(\sum_{j=1}^n \norm{a_{ij}}\biggr)^2 \le \biggl(\sum_{i=1}^n\sum_{j=1}^n \norm{a_{ij}}\biggr)^2 \]
setiap kali $\vc a \in \M nA$, $\vc v \in H^n$, dan $\norm{\vc v} \le 1$.    \ns
\end{proof}


























\section{Semigrup Proyeksi}

\begin{notn} Jika $A$ suatu aljabar-$C^*$ dan $n \in \N$, kita menetapkan
 \index{P@$\fml P_n(A)$, $\fml P_\infty(A)$ (proyeksi dalam aljabar matriks)}%
  \begin{align*}
       \fml P_n(A) &= \fml P(\M nA) \quad\text{ dan} \\
       \fml P_\infty(A) &= \bigcup_{n=1}^\infty \fml P_n(A).
  \end{align*}
\end{notn}

Sekarang kita memperluas ekuivalensi Murray--von Neumann ke matriks-matriks yang ukurannya berbeda.

\begin{notn} Jika $A$ suatu aljabar-$C^*$, misalkan
 \index{Ma@$\M {n,m}A$ (matriks dengan entri dari~$A$)}%
$\M {n,m}A$ menyatakan himpunan matriks $n \times m$ berentri dalam~$A$.
\end{notn}

Selanjutnya kita memperluas ekuivalensi Murray--von Neumann ke matriks proyeksi yang ukurannya tidak harus sama.

\begin{defn}\label{0060316fa} Jika $A$ suatu aljabar-$C^*$, $p \in \fml P_m(A)$, dan $q \in \fml P_n(A)$, kita menetapkan
 \index{ekuivalensi Murray--von Neumann}%
 \index{ekuivalensi!Murray--von Neumann}%
 \index{<binrelequiv@$p \sim q$ (ekuivalensi Murray--von Neumann)}%
$p \sim q$ jika terdapat $v \in \M {n,m}A$ sedemikian sehingga
  \[ v^*v = p \qquad\text{ dan } \qquad  vv^* = q. \]
Kita akan memperluas penggunaan nama \emph{ekuivalensi Murray--von Neumann} untuk relasi baru ini pada~$\fml P_\infty(A)$.
\end{defn}

\begin{prop}\label{0060317} Relasi $\sim$ yang didefinisikan di atas merupakan relasi
ekuivalensi pada~$\fml P_\infty(A)$.
\end{prop}

\begin{defn} Untuk setiap aljabar-$C^*$ $A$, kita mendefinisikan operasi biner $\oplus$ pada $\fml
P_\infty(A)$ dengan
   \[ p \oplus q = \diag(p,q)\,. \]
Jadi, jika $p \in \fml P_m(A)$ dan $q \in \fml P_n(A)$, maka $p \oplus q \in \fml
P_{m+n}(A)$.
\end{defn}

Dalam proposisi berikut, $\vc 0_n$ adalah identitas aditif dalam~$\fml P_n(A)$.

\begin{prop}\label{0060324} Misalkan $A$ suatu aljabar-$C^*$ dan $p \in \fml P_\infty(A)$.
 \index{<sum@$p \oplus q$ (operasi biner dalam $\fml P_\infty(A)$)}%
Maka $p \sim p \oplus \vc 0_n$ untuk setiap $n \in \N$.
\end{prop}

\begin{prop}\label{0060327} Misalkan $A$ suatu aljabar-$C^*$ dan $p$, $p'$, $q$, $q' \in
\fml P_\infty(A)$. Jika $p \sim p'$ dan $q \sim q'$, maka $p \oplus q \sim p' \oplus q'$.
\end{prop}

\begin{prop}\label{0060331} Misalkan $A$ suatu aljabar-$C^*$ dan $p$, $q \in \fml P_\infty(A)$.
Maka $p \oplus q \sim q \oplus p$.
\end{prop}

\begin{prop}\label{0060334} Misalkan $A$ suatu aljabar-$C^*$ dan $p$, $q \in \fml P_n(A)$ untuk
suatu $n$. Jika $p \perp q$, maka $p + q$ suatu proyeksi dalam $\M nA$ dan $p + q \sim p \oplus q$.
\end{prop}

\begin{prop}\label{0060337} Jika $A$ suatu aljabar-$C^*$, maka operasi $\oplus$ pada $\fml P_\infty(A)$ bersifat
asosiatif secara ketat dan komutatif hanya hingga ekuivalensi Murray--von Neumann. % SOURCE-CORRECTION: CH17-C006
\end{prop}

\begin{notn}\label{0060411} Jika $A$ suatu aljabar-$C^*$ dan $p \in \fml P_\infty(A)$, misalkan
 \index{<bracsqr@$\sbsb{[p]}{\fml D}$ (kelas ekuivalensi $\sim$ dari $p$)}%
$\sbsb{[p\,]}{\fml D}$ adalah kelas ekuivalensi yang memuat $p$ dan ditentukan oleh
relasi ekuivalensi~$\sim$. Tetapkan pula
 \index{D@$\fml D(A)$ (himpunan kelas ekuivalensi proyeksi)}%
$\fml D(A) := \{\,\sbsb{[p\,]}{\fml D}\colon p \in \fml P_\infty(A)\,\}$.
\end{notn}

\begin{defn}\label{0060414} Misalkan $A$ suatu aljabar-$C^*$. Definisikan operasi biner $+$ pada
$\fml D(A)$ dengan
 \index{<binop@$\sbsb{[p\,]}{\fml D} + \sbsb{[q\,]}{\fml D}$ (penjumlahan dalam $\fml D(A)$)}%
   \[ \sbsb{[p\,]}{\fml D} + \sbsb{[q\,]}{\fml D} := \sbsb{[p \oplus q\,]}{\fml D} \]
dengan $p$, $q \in \fml P_\infty(A)$.
\end{defn}

\begin{prop}\label{0060417} Operasi $+$ yang didefinisikan dalam~\ref{0060414} terdefinisi dengan baik
dan menjadikan $\fml D(A)$ suatu semigrup komutatif.
\end{prop}

\begin{exam}\label{0060421} Semigrup $\fml D(\C)$ isomorfik dengan semigrup aditif bilangan bulat tak negatif; yaitu, % SOURCE-CORRECTION: CH17-C007
    \[ \fml D(\C) \cong \Z^+ = \{0,1,2,\dots\}\,.  \]
\end{exam}

\begin{proof}[\emph{Petunjuk untuk bukti}] Gunakan contoh~\ref{0060247fa}. \ns
\end{proof}

\begin{exam}\label{0060424} Jika $H$ suatu ruang Hilbert terpisahkan berdimensi tak hingga, maka % SOURCE-CORRECTION: CH17-C008
   \[ \fml D(\ofml B(H)) \cong \Z^+ \cup \{\infty\}. \]
(Gunakan penjumlahan biasa pada $\Z^+$ dan tetapkan $n + \infty = \infty + n = \infty$ setiap kali $n
\in \Z^+ \cup \{\infty\}$.)
\end{exam}

\begin{exam}\label{0060427} Untuk aljabar-$C^*$ $\C \oplus \C$, kita mempunyai
   \[ \fml D(\C \oplus \C) \cong \Z^+ \oplus \Z^+\,. \]
\end{exam}


























\section{Konstruksi Grothendieck}
Dalam konstruksi medan bilangan real, kita dapat menggunakan teknik yang sama untuk beralih dari (semigrup aditif) bilangan asli ke
(grup aditif semua) bilangan bulat seperti yang kita gunakan untuk beralih dari (semigrup multiplikatif) bilangan bulat tak nol ke (grup multiplikatif
semua) bilangan rasional tak nol. Teknik ini disebut \emph{konstruksi Grothendieck}.

\begin{defn}\label{0062111} Misalkan $(S,+)$ suatu semigrup komutatif. Definisikan relasi
 \index{<binrelequiv@$(a,b) \sim (c,d)$ (ekuivalensi Grothendieck)}%
$\sim$ pada $S \times S$ dengan
   \[ (a,b) \sim (c,d) \text{ jika terdapat $k \in S$ sedemikian sehingga } a + d + k = b + c + k. \]
\end{defn}

\begin{prop}\label{0062114} Relasi $\sim$ yang didefinisikan di atas merupakan relasi ekuivalensi.
\end{prop}

\begin{notn} Untuk relasi ekuivalensi $\sim$ yang didefinisikan dalam~\ref{0062111}, kelas ekuivalensi
yang memuat pasangan $(a,b)$ akan dinotasikan dengan % SOURCE-CORRECTION: CH17-C009
 \index{<bracpnt@$\langle a,b \rangle$ (pasangan dalam konstruksi Grothendieck)}%
$\langle a,b \rangle$, bukan dengan~$[(a,b)]$.
\end{notn}

\begin{defn} Misalkan $(S,+)$ suatu semigrup komutatif. Pada $G(S)$, definisikan operasi biner
 \index{penjumlahan!pada grup Grothendieck}%
(yang juga dinotasikan dengan $+$) sebagai berikut:
  \[ \langle a,b \rangle + \langle c,d \rangle := \langle a + c,b + d\rangle\,. \]
\end{defn}

\begin{prop}\label{0062119} Operasi $+$ yang didefinisikan di atas terdefinisi dengan baik, dan di bawah
operasi ini $G(S)$ menjadi grup Abelian. % SOURCE-CORRECTION: CH17-C010
\end{prop}

Grup Abelian $(G(S),+)$ disebut
 \index{G@$G(S)$ (grup Grothendieck dari $S$)}%
 \index{Grothendieck!grup}%
 \index{grup!Grothendieck}%
\df{grup Grothendieck} dari~$S$.

\begin{prop}\label{0062124} Untuk suatu semigrup $S$ dan sembarang $a \in S$, definisikan
pemetaan
 \index{gamma@$\sbsb{\gamma}S$ (pemetaan Grothendieck)}%
   \[ \sbsb{\gamma}S \colon S \sto G(S) \colon s \mapsto \langle s + a, a \rangle\,. \]
Pemetaan $\sbsb{\gamma}S$, yang disebut
 \index{Grothendieck!pemetaan}%
\df{pemetaan Grothendieck}, terdefinisi dengan baik dan merupakan homomorfisme semigrup.
\end{prop}

Pernyataan bahwa pemetaan Grothendieck terdefinisi dengan baik berarti bahwa definisinya tidak bergantung
pada pilihan~$a$. Kita sering hanya menulis $\gamma$ untuk~$\sbsb{\gamma}S$.

\begin{exam}\label{0062126} Baik $\N = \{1,2,3,\dots\}$ maupun $\Z^+ = \{0,1,2,\dots\}$ merupakan
semigrup komutatif di bawah penjumlahan. Keduanya menghasilkan grup Grothendieck yang sama,
   \[G(\N) = G(\Z^+) = \Z\,. \]
\end{exam}

Tidak ada yang menjamin bahwa grup Grothendieck dari suatu semigrup sembarang akan sangat menarik.

\begin{exam}\label{0062127} Misalkan $S$ semigrup aditif komutatif $\Z^+ \cup \{\infty\}$. Maka $G(S) = \{\vc 0\}$.
\end{exam}

\begin{exam}\label{0062128} Misalkan $\Z_0$ semigrup multiplikatif (komutatif) bilangan bulat tak nol.
Maka $G(\Z_0) = \Q_0$, yaitu grup multiplikatif Abelian dari bilangan rasional tak
nol.
\end{exam}

\begin{prop}\label{0062141} Jika $S$ suatu semigrup komutatif, maka
   \[ G(S) = \{\gamma(s) - \gamma(t)\colon s,t \in S\}\,. \]
\end{prop}

\begin{prop}\label{0062142} Jika $r$, $s \in S$, dengan $S$ suatu semigrup komutatif, maka $\gamma(r) = \gamma(s)$ jika dan hanya jika
terdapat $t \in S$ sedemikian sehingga $r + t = s + t$.
\end{prop}

\begin{defn} Suatu semigrup komutatif $S$ memiliki
 \index{sifat pembatalan}%
\df{sifat pembatalan} jika setiap kali $r$, $s$, $t \in S$ memenuhi $r + t = s + t$, berlaku $r = s$.
\end{defn}

\begin{cor}\label{0062144} Misalkan $S$ suatu semigrup komutatif. Pemetaan Grothendieck $\sbsb\gamma S\colon S \sto G(S)$ injektif
jika dan hanya jika $S$ memiliki sifat pembatalan.
\end{cor}

Proposisi berikut menyatakan sifat universal grup Grothendieck.

\begin{prop}\label{0062151} Misalkan $S$ suatu semigrup (aditif) komutatif dan $G(S)$ grup Grothendieck-nya.
 \index{universal!sifat!pemetaan Grothendieck}%
 \index{Grothendieck!pemetaan!sifat universal}%
Jika $H$ suatu grup Abelian dan $\phi\colon S \sto H$ suatu pemetaan aditif, maka terdapat
homomorfisme grup tunggal $\psi\colon G(S) \sto H$ sedemikian sehingga diagram berikut
komutatif.
 \begin{equation}\label{0062151i}
   \xy
     \qtriangle/{>}`{>}`{-->}/[S`\abs{G(S)}`\abs{H};\gamma`\phi`\abs{\psi}]
     \morphism(1100,500)|r|/{-->}/<0,-500>[G(S)`H;\psi]
   \endxy
 \end{equation}
\end{prop}

Dalam diagram sebelumnya, $\abs{G(S)}$ dan $\abs{H}$ hanyalah $G(S)$ dan $H$ yang dipandang sebagai
semigrup, sedangkan $\abs{\psi}$ adalah homomorfisme semigrup yang bersesuaian. Dengan kata lain,
funktor pelupa $\abs{\hphantom{G}}$ ``melupakan'' hanya identitas dan invers, tetapi tidak
operasi penjumlahan. Jadi, segitiga di sebelah kiri merupakan diagram komutatif dalam
kategori semigrup dan homomorfisme semigrup.

\begin{prop}\label{0062154} Misalkan $\phi\colon S \sto T$ homomorfisme semigrup
komutatif. Maka pemetaan $\sbsb{\gamma}T \circ \phi\colon S \sto G(T)$ aditif. Berdasarkan
proposisi~\ref{0062151}, terdapat homomorfisme grup tunggal $G(\phi)\colon G(S)
\sto G(T)$ sedemikian sehingga diagram berikut komutatif.
   \[ \xy
          \Square[S`T`G(S)`G(T);\phi`\sbsb{\gamma}S`\sbsb{\gamma}T`G(\phi)]
      \endxy \]
\end{prop}

\begin{prop}\label{0062155} Pasangan pemetaan $S \mapsto G(S)$, yang memetakan semigrup komutatif ke
 \index{funktor!Grothendieck}%
 \index{Grothendieck!konstruksi!sifat funktorial}%
grup Grothendieck yang bersesuaian, dan $\phi \mapsto G(\phi)$, yang memetakan homomorfisme semigrup ke
homomorfisme-homomorfisme grup (sebagaimana didefinisikan dalam~\ref{0062154}), merupakan funktor kovarian dari kategori % SOURCE-CORRECTION: CH17-C011
semigrup komutatif dan homomorfisme semigrup ke kategori grup Abelian
dan homomorfisme grup.
\end{prop}

Salah satu keuntungan kecil dari penyertaan funktor pelupa yang agak bertele-tele dalam
diagram~\eqref{0062151i} adalah bahwa hal itu memungkinkan kita memandang pemetaan Grothendieck
$\gamma\colon S \mapsto \sbsb{\gamma}S$ sebagai transformasi alami antarfungtor.

\begin{cor}[Kealamian Pemetaan Grothendieck]\label{0062157} Misalkan $\abs{\hphantom{G}}$ funktor pelupa
pada grup Abelian yang ``melupakan'' identitas dan invers, tetapi tidak
operasi grup, seperti dalam~\ref{0062151}. Maka $\abs{G(\hphantom{S})}$ merupakan funktor kovarian dari
kategori semigrup komutatif dan homomorfisme semigrup ke dirinya sendiri.
 \index{Grothendieck!pemetaan!kealamian}%
 \index{alami!transformasi!pemetaan Grothendieck sebagai suatu}%
Selanjutnya, pemetaan Grothendieck $\gamma\colon S \mapsto \sbsb{\gamma}S$ merupakan transformasi
alami dari funktor identitas ke funktor~$\abs{G(\hphantom{S})}$.
  \[ \xy
          \Square[S`T`\abs{G(S)}`\abs{G(T)};\phi`\sbsb{\gamma}S`\sbsb{\gamma}T`\abs{G(\phi)}]
      \endxy \]
\end{cor}





\section{Grup $\mathbf{\emph{K}_0}$ untuk Aljabar-$C^*$ Beridentitas}

\begin{defn} Misalkan $A$ aljabar-$C^*$ beridentitas. Misalkan
 \index{K0@$K_0(A)$!untuk $A$ beridentitas}%
$K_0(A) := G(\fml D(A))$, grup Grothendieck dari semigrup $\fml D(A)$ yang didefinisikan dalam~\ref{0060411} dan~\ref{0060414},
dan definisikan
 \index{<bracsqr@$[p\,]$ (citra $[p\,]_{\fml D}$ di bawah pemetaan Grothendieck)}%
   \[ [\hphantom{P}]\colon \fml P_\infty(A) \sto K_0(A)
           \colon p \mapsto \sbsb{\gamma}{\fml D(A)}\bigl([p\,]_{\fml D}\bigr)\,. \]
\end{defn}

Proposisi berikut merupakan konsekuensi langsung dari definisi ini.

\begin{prop}\label{0062188fa} Misalkan $A$ aljabar-$C^*$ beridentitas dan $p$, $q \in \fml P_\infty(A)$. Jika $p$ dan $q$
ekuivalen Murray--von Neumann, maka $[p\,] = [q\,]$ dalam~$K_0(A)$.
\end{prop}

\begin{defn} Misalkan $A$ aljabar-$C^*$ beridentitas dan $p$, $q \in \fml P_\infty(A)$. Kita mengatakan bahwa $p$
 \index{ekuivalensi!stabil}%
 \index{ekuivalensi stabil}%
\df{ekuivalen secara stabil} dengan $q$ dan menulis
 \index{<binrelequivst@$p \sim_{st} q$ (ekuivalensi stabil)}%
$p \sim_{st} q$ jika terdapat proyeksi $r \in \fml P_\infty(A)$ sedemikian sehingga $p \oplus r \sim q \oplus r$.
\end{defn}

\begin{prop} Ekuivalensi stabil $\sim_{st}$ merupakan relasi ekuivalensi pada $\fml P_\infty(A)$.
\end{prop}

\begin{prop} Misalkan $A$ aljabar-$C^*$ beridentitas dan $p$, $q \in \fml P_\infty(A)$. Maka $p
\sim_{st} q$ jika dan hanya jika $p \oplus \vc 1_n \sim q \oplus \vc 1_n$ untuk suatu $n \in \N$.
\end{prop}

Di sini, tentu saja, $\vc 1_n$ adalah identitas perkalian dalam $\M nA$.

\begin{prop}[Gambaran Standar $K_0(A)$ ketika $A$ beridentitas]\label{0062224} Jika $A$ suatu
aljabar-$C^*$ beridentitas, maka
 \index{standar!gambaran $K_0(A)$!ketika $A$ beridentitas}%
 \index{gambaran (\seeonly{gambaran standar}))}%
   \begin{align*}
         K_0(A) &= \{ [p\,] - [q\,]\colon p,q \in \fml P_\infty(A)\}  \\
                &= \{ [p\,] - [q\,]\colon  p,q \in \fml P_n(A) \text{ untuk suatu $n \in \N$}\}\,.
   \end{align*}
\end{prop}

\begin{prop} Misalkan $A$ aljabar-$C^*$ beridentitas dan $p$, $q \in \fml P_\infty(A)$. Maka $[p
\oplus q\,] = [p\,] + [q\,]$.
\end{prop}

\begin{prop} Misalkan $A$ aljabar-$C^*$ beridentitas dan $p$, $q \in \fml P_n(A)$. Jika $p \sim_h q$
dalam $\fml P_n(A)$, maka $[p\,] = [q\,]$.
\end{prop}

\begin{prop} Misalkan $A$ aljabar-$C^*$ beridentitas dan $p$, $q \in \fml P_n(A)$. Jika $p \perp q$
dalam $\fml P_n(A)$, maka $p + q \in \fml P_n(A)$ dan $[p + q\,] = [p\,] + [q\,]$.
\end{prop}

\begin{prop} Misalkan $A$ aljabar-$C^*$ beridentitas dan $p$, $q \in \fml P_\infty(A)$. Maka
$[p\,] = [q\,]$ jika dan hanya jika $p \sim_{st} q$.
\end{prop}

 Proposisi berikut menetapkan suatu sifat universal dari $K_0(A)$ ketika $A$ beridentitas. Di dalamnya, funktor pelupa $\abs{\hphantom{G}}$
 adalah funktor yang dijelaskan dalam proposisi~\ref{0062151}.

\begin{prop}\label{0062244} Misalkan $A$ aljabar-$C^*$ beridentitas, $G$ grup Abelian, dan $\nu\colon \fml P_\infty(A) \sto \abs G$
homomorfisme semigrup yang memenuhi
 \index{universal!sifat!dari $K_0$ (kasus beridentitas)}%
 \index{K0@$K_0$!sifat universal dari}%
 \begin{enumerate}[label=\rm{(\alph*)}]
    \item $\nu(\vc 0_A) = \vc 0_G$ dan
    \item jika $p \sim_h q$ dalam $\fml P_n(A)$ untuk suatu $n$, maka $\nu(p) = \nu(q)$.
 \end{enumerate}
Maka terdapat homomorfisme grup tunggal $\wt\nu\colon K_0(A) \sto G$ sedemikian sehingga diagram berikut komutatif.
   \[\xy
      \qtriangle/{>}`{>}`{-->}/<1000,700>[\fml P_\infty(A)`\abs{K_0(A)}`\abs G;{[\hphantom{P}]}`\nu`\abs{\wt\nu}]
      \morphism(1600,700)|r|/{-->}/<0,-700>[K_0(A)`G;\wt\nu]
     \endxy\]
\end{prop}

\begin{proof}[\emph{Petunjuk untuk bukti}] Misalkan $\tau\colon \fml D(A) \sto \abs G\colon \sbsb{[p\,]}{\fml D} \mapsto \nu(p)$. % SOURCE-CORRECTION: CH17-C026
Periksa bahwa $\tau$ terdefinisi dengan baik dan merupakan homomorfisme semigrup. Kemudian gunakan proposisi~\ref{0062151}.  \ns
\end{proof}

\begin{defn} Homomorfisme-$*\,$ $\phi\colon A \sto B$ di antara aljabar-$C^*$ diperluas, untuk
setiap $n \in \N$, menjadi homomorfisme-$*\,$ $\phi\colon \M nA \sto \M nB$ dan juga (karena homomorfisme-$*\,$ membawa proyeksi
ke proyeksi) memiliki pembatasan berupa pemetaan $\phi$ dari $\fml P_\infty(A)$ ke $\fml P_\infty(B)$. % SOURCE-CORRECTION: CH17-C012
Untuk homomorfisme-$*\,$ $\phi$ semacam itu, definisikan
   \[ \nu\colon \fml P_\infty(A) \sto K_0(B)\colon p \mapsto [\phi(p)]\,. \]
Maka $\nu$ merupakan homomorfisme semigrup yang memenuhi syarat (a) dan (b) dalam proposisi~\ref{0062244}, sehingga
terdapat homomorfisme grup tunggal
 \index{K0@$K_0(\phi)$!kasus beridentitas}%
$K_0(\phi)\colon K_0(A) \sto K_0(B)$ sedemikian sehingga $K_0(\phi)\bigl([p\,]\bigr) = \nu(p)$ untuk setiap $p \in \fml P_\infty(A)$.
\end{defn}

\begin{prop} Pasangan pemetaan $A \mapsto K_0(A)$, $\phi \mapsto K_0(\phi)$ merupakan funktor kovarian
dari kategori aljabar-$C^*$ beridentitas dan homomorfisme-$*\,$ ke kategori grup Abelian dan
homomorfisme grup. Lebih lanjut, untuk semua homomorfisme-$*\,$ $\phi\colon A \sto B$ di antara aljabar-$C^*$ beridentitas, diagram
berikut komutatif.
  \[\xy
         \Square[\fml P_\infty(A)`\fml
         P_\infty(B)`K_0(A)`K_0(B);\phi`{[\hphantom{P}]}`{[\hphantom{P}]}`K_0(\phi)]
  \endxy\]
\end{prop}

\begin{notn} Untuk aljabar-$C^*$ $A$ dan $B$, misalkan
 \index{hom@$\Hom(A,B)$ (himpunan homomorfisme-$*\,$ di antara aljabar-$C^*$)}%
$\Hom(A,B)$ adalah keluarga semua homomorfisme-$*\,$ dari $A$ ke~$B$.
\end{notn}

\begin{defn} Misalkan $A$ dan $B$ aljabar-$C^*$ serta $a \in A$. Untuk $\phi$, $\psi \in \Hom(A,B)$,
misalkan
   \[ d_a(\phi,\psi) = \norm{\phi(a) - \psi(a)}\,. \]
Maka $d_a$ merupakan pseudometrik pada $\Hom (A,B)$. Topologi yang dibangkitkan oleh keluarga $\{d_a\colon a \in A\}$ adalah
 \index{topologi norma-titik}%
 \index{topologi!norma-titik}%
\df{topologi norma-titik} pada~$\Hom(A,B)$. (Untuk setiap $a \in A$, misalkan $\sfml B_a$ adalah keluarga bola terbuka dalam $\Hom(A,B)$ yang dibangkitkan oleh
pseudometrik~$d_a$. Keluarga $\bigcup\{\sfml B_a\colon a \in A\}$ merupakan subbasis bagi topologi norma-titik.)
\end{defn}

\begin{defn} Misalkan $A$ dan $B$ aljabar-$C^*$. Kita mengatakan bahwa homomorfisme-$*\,$ $\phi$, $\psi\colon A \sto B$
 \index{homotop!homomorfisme-$*\,$}%
 \index{<star@homomorfisme-$*\,$!homotopi dari}%
 \index{<binrelequivh@$\phi\sim_h\psi$ (homotopi homomorfisme-$*\,$)}%
\df{homotop}, dan menulis $\phi \sim_h \psi$, jika terdapat suatu fungsi
   \[ c\colon [0,1] \times A \sto B\colon (t,a) \mapsto c_t(a) \]
sedemikian sehingga
  \begin{enumerate}[label=\rm{(\alph*)}]
     \item untuk setiap $t \in [0,1]$, pemetaan $c_t\colon A \sto B \colon a \mapsto c_t(a)$ merupakan homomorfisme-$*\,$,
     \item untuk setiap $a \in A$, pemetaan $c(a)\colon [0,1] \sto B\colon t \mapsto c_t(a)$ merupakan lintasan (kontinu) dalam~$B$,
     \item $c_0 = \phi$, dan
     \item $c_1 = \psi$.
  \end{enumerate}
\end{defn}

\begin{prop} Dua homomorfisme-$*\,$ $\phi_0$, $\phi_1\colon A \sto B$ di antara aljabar-$C^*$
homotop jika dan hanya jika terdapat lintasan kontinu norma-titik dari $\phi_0$ ke $\phi_1$ dalam~$\Hom(A,B)$.
\end{prop}

\begin{defn} Kita mengatakan bahwa aljabar-$C^*$ $A$ dan $B$
 \index{homotop!ekuivalensi!homomorfisme-$*\,$}%
 \index{ekuivalensi!homotopi!homomorfisme-$*\,$}%
\df{ekuivalen secara homotopi} jika terdapat homomorfisme-$*\,$ $\phi\colon A \sto B$ dan $\psi\colon B \sto A$ sedemikian sehingga $\psi
\circ \phi \sim_h \id A$ dan $\phi \circ \psi \sim_h \id B$. Suatu aljabar-$C^*$ disebut
 \index{kontraktibel!aljabar-$C^*$}%
\df{kontraktibel} jika ekuivalen secara homotopi dengan~$\{\vc 0\}$.
\end{defn}

\begin{prop}\label{0062327} Misalkan $A$ dan $B$ aljabar-$C^*$ beridentitas. Jika $\phi \sim_h \psi$
dalam $\Hom(A,B)$, maka $K_0(\phi) = K_0(\psi)$.
\end{prop}

%\begin{proof} Lihat \cite{RordamLL:2000}, proposisi 3.2.6.  \ns \end{proof}

\begin{prop}\label{0062331} Jika aljabar-$C^*$ beridentitas $A$ dan $B$
 \index{homotopi!invariansi $K_0$}%
 \index{K0@$K_0$!invariansi homotopi dari}%
ekuivalen secara homotopi, maka $K_0(A) \cong K_0(B)$.
\end{prop}

\begin{prop}\label{0062334} Jika $A$ suatu aljabar-$C^*$ beridentitas, maka barisan eksak terbelah
   \begin{equation}\label{0062334i}
      \vc 0 \to A \to^\iota \wt A \two/->`<-/^Q_\psi \C \to \vc 0
   \end{equation}
\emph{(lihat \ref{0015213}, Kasus 1, butir (7)\,)} menginduksi barisan eksak terbelah lainnya
   \begin{equation}\label{0062334ii}
      \vc 0 \to K_0(A) \to^{K_0(\iota)} K_0(\wt A) \two/->`<-/^{K_0(Q)}_{K_0(\psi)} K_0(\C) \to \vc 0 \,.
   \end{equation}
\end{prop}

\begin{proof}[\emph{Petunjuk untuk bukti}] Definisikan dua fungsi tambahan
   \[ \mu\colon \wt A \sto A\colon a + \lambda \vc j \mapsto a \]
dan
   \[ \psi' \colon \C \sto \wt A \colon \lambda \mapsto \lambda \vc j\,. \]
Kemudian periksa bahwa
   \begin{enumerate}[label=\rm{(\alph*)}]
     \item $\mu \circ \iota = \id A$,
     \item $\iota \circ \mu + \psi' \circ Q = \id{\wt A}$,
     \item $Q \circ \iota = \vc 0$, dan
     \item $Q \circ \psi' = \id{\C}$. % SOURCE-CORRECTION: CH17-C013
   \end{enumerate}   \ns
\end{proof}

\begin{exam} Untuk setiap $n \in \N$,
 \index{K0@$K_0(A)$!adalah $\Z$ ketika $A = \mathbf M_n$}%
$K_0(\mathbf M_n) \cong \Z$.
\end{exam}

\begin{exam} Jika $H$ ruang Hilbert terpisahkan berdimensi tak hingga,
 \index{K0@$K_0(A)$!adalah $\vc 0$ ketika $A = \ofml B(H)$}%
maka $K_0(\ofml B(H)) \cong \vc 0$.
\end{exam}

\begin{defn} Ingat bahwa suatu ruang topologis $X$ disebut
 \index{kontraktibel!ruang topologis}%
\df{kontraktibel} jika terdapat titik $a$ dalam ruang tersebut dan fungsi kontinu $f\colon [0,1] \times X \sto X$ sedemikian sehingga $f(1,x) = x$
dan $f(0,x) = a$ untuk setiap $x \in X$.
\end{defn}

\begin{exam} Jika $X$ ruang Hausdorff kompak yang kontraktibel, maka
 \index{K0@$K_0(A)$!adalah $\Z$ ketika $A = \fml C(X)$, $X$ kontraktibel}%
$K_0(\fml C(X)) \cong \Z$.
\end{exam}


















\section{$\mathbf{\emph{K}_0}(A)$---Kasus Tak Beridentitas}
\begin{defn} Misalkan $A$ aljabar-$C^*$ tak beridentitas. Ingat bahwa barisan eksak terbelah
 ~\eqref{0062334i} bagi unitalisasi $A$ menginduksi barisan eksak terbelah
~\eqref{0062334ii} di antara grup-grup $K_0$ yang bersesuaian.
 \index{K0@$K_0(A)$!untuk $A$ tak beridentitas}%
Definisikan
   \[ \pi:=Q,\qquad K_0(A) = \ker(K_0(\pi))\,. \] % SOURCE-CORRECTION: CH17-C014
\end{defn}

\begin{prop}\label{0062413} Untuk aljabar-$C^*$ tak beridentitas $A$, pemetaan $[\hphantom{P}] \colon \fml
P_\infty(A) \sto K_0(\wt A)$ dapat dipandang sebagai pemetaan dari $\fml P_\infty(A)$ ke dalam~$K_0(A)$.
\end{prop}

\begin{prop}\label{0062416} Untuk aljabar-$C^*$ beridentitas maupun tak beridentitas, barisan
  \[ \vc 0 \to K_0(A) \to K_0(\wt A) \to K_0(\C) \to \vc 0 \]
bersifat eksak.
\end{prop}

\begin{prop}\label{0062418} Untuk aljabar-$C^*$ beridentitas maupun tak beridentitas, grup $K_0(A)$ (isomorfik
dengan) $\ker(K_0(\pi))$.
\end{prop}

\begin{prop}\label{0062424} Jika $\phi\colon A \sto B$ merupakan homomorfisme-$*\,$ di antara aljabar-$C^*$, maka
terdapat homomorfisme grup tunggal % SOURCE-CORRECTION: CH17-C015
 \index{K0@$K_0(\phi)$!kasus tak beridentitas}%
$K_0(\phi)$ yang membuat diagram berikut komutatif.
   \[\xy
     \hSquares/>`>`.>`>`=`>`>/[K_0(A)`K_0(\wt A)`K_0(\C)`K_0(B)`K_0(\wt
     B)`K_0(\C);`K_0(\pi_A)`K_0(\phi)`K_0(\wt\phi)```K_0(\pi_B)]
   \endxy\]
\end{prop}

\begin{prop}\label{0062511} Pasangan pemetaan $A \mapsto K_0(A)$, $\phi \mapsto K_0(\phi)$ merupakan
 \index{funktor!$K_0$}%
 \index{K0@$K_0$!sebagai funktor kovarian}%
funktor kovarian dari kategori $\cat{CSA}$ aljabar-$C^*$ dan homomorfisme-$*\,$ ke
kategori grup Abelian dan homomorfisme grup.
\end{prop}

Dalam proposisi~\ref{0062327} dan~\ref{0062331}, kita menyatakan invariansi homotopi dari
funktor $K_0$ untuk aljabar-$C^*$ beridentitas. Sekarang kita memperluas hasil tersebut ke sebarang
aljabar-$C^*$.

\begin{prop}\label{0062521} Misalkan $A$ dan $B$ aljabar-$C^*$. Jika $\phi \sim_h \psi$
dalam $\Hom(A,B)$, maka $K_0(\phi) = K_0(\psi)$.
\end{prop}

%\begin{proof} Lihat \cite{RordamLL:2000}, proposisi 4.1.4.  \ns \end{proof}

\begin{prop}\label{0062524} Jika aljabar-$C^*$ $A$ dan $B$ ekuivalen secara homotopi,
maka $K_0(A) \cong K_0(B)$.
\end{prop}

%\begin{proof} Lihat \cite{RordamLL:2000}, proposisi 4.1.4.  \ns \end{proof}

\begin{defn} Misalkan $\pi:=Q$ dan $\lambda:=\psi$ adalah homomorfisme-$*\,$ yang berturut-turut berperan sebagai pemetaan hasil bagi dan penampang kanonik dalam barisan eksak
terbelah~\eqref{0062334i} bagi unitalisasi suatu aljabar-$C^*$~$A$. Definisikan % SOURCE-CORRECTION: CH17-C016
 \index{skalar!pemetaan}%
\df{pemetaan skalar} $s\colon \wt A \sto \wt A$ untuk $\wt A$ dengan $s :=  \lambda \circ \pi$. Setiap
anggota $\wt A$ dapat ditulis dalam bentuk $a + \alpha \vc 1_{\wt A}$ untuk suatu $a \in A$ dan
$\alpha \in \C$. Perhatikan bahwa $s(a + \alpha \vc 1_{\wt A}) = \alpha \vc 1_{\wt A}$ dan bahwa $x
- s(x) \in A$ untuk setiap $x \in \wt A$. Untuk setiap bilangan asli $n$, pemetaan skalar menginduksi
pemetaan yang bersesuaian $s = s_n\colon \M n{\wt A} \sto \M n{\wt A}$. Suatu elemen $x \in \M n{\wt
A}$ merupakan
 \index{skalar!elemen}%
\df{elemen skalar} dari $\M n{\wt A}$ jika $s(x) = x$.
\end{defn}

\begin{prop}[Gambaran Standar $K_0(A)$ untuk sebarang $A$]\label{0062544} Jika $A$ suatu
aljabar-$C^*$,
 \index{standar!gambaran $K_0(A)$!untuk sebarang $A$}%
 \index{gambaran (\seeonly{gambaran standar}))}%
maka
     \[ K_0(A) = \{\,[p\,] - [s(p)]\colon p \in \fml P_\infty(\wt A)\}\,. \] % SOURCE-CORRECTION: CH17-C017
\end{prop}























\section{Sifat Keeksakan dan Stabilitas Funktor $K_0$}
\begin{defn} Funktor kovarian $F$ dari kategori $\cat A$ ke kategori $\cat B$ disebut
 \index{eksak terbelah!funktor}%
 \index{funktor!eksak terbelah}%
\df{eksak terbelah} jika membawa barisan eksak terbelah ke barisan eksak terbelah. Funktor tersebut disebut
 \index{funktor eksak separuh}%
 \index{funktor!eksak separuh}%
\df{eksak separuh} jika setiap kali barisan
   \[ \vc 0 \to A_1 \to^j A_2 \to^k A_3 \to \vc 0 \]
eksak dalam $\cat A$, maka
   \[ F(A_1) \to^{F(j)} F(A_2) \to^{F(k)} F(A_3) \]
eksak dalam~$\cat B$.
\end{defn}

\begin{prop}\label{0062664} Funktor $K_0$ eksak separuh.
\end{prop}

%\begin{proof} Lihat~\cite{RordamLL:2000}, proposisi 4.3.2.  \ns \end{proof}

\begin{prop}\label{0062667} Funktor $K_0$ eksak terbelah.
\end{prop}

%\begin{proof} Lihat~\cite{RordamLL:2000}, proposisi 4.3.3.  \ns \end{proof}

\begin{prop}\label{0062671} Funktor $K_0$ mempertahankan jumlah langsung. Artinya, jika $A$ dan $B$
aljabar-$C^*$, maka $K_0(A \oplus B) = K_0(A) \oplus K_0(B)$.
\end{prop}

\begin{exam}\label{0062674} Jika $A$ aljabar-$C^*$, maka $K_0(\wt A) = K_0(A) \oplus \Z$.
\end{exam}

Meskipun eksak terbelah sekaligus eksak separuh, funktor $K_0$ tidak eksak. Masing-masing dari dua
contoh berikut cukup untuk menunjukkan hal ini.

\begin{exam}\label{0062677} Barisan
   \[ \vc 0 \to \fml C_0\bigl((0,1)\bigr) \to^\iota \fml C\bigl([0,1]\bigr)
                               \to^\psi \C \oplus \C \to \vc 0 \]
dengan $\psi(f) = \bigl(f(0),f(1)\bigr)$ jelas eksak; tetapi $K_0(\psi)$ tidak surjektif.
\end{exam}

\begin{exam}\label{0062679} Jika $H$ ruang Hilbert berdimensi tak hingga, barisan eksak % SOURCE-CORRECTION: CH17-C018
  \[ \vc 0 \to \ofml K(H) \to^\iota \ofml B(H) \to^\pi \ofml Q(H) \to \vc 0 \]
yang terkait dengan aljabar Calkin $\ofml Q(H)$ bersifat eksak, tetapi $K_0(\iota)$ tidak injektif.
(Contoh ini memerlukan fakta yang belum kita turunkan: $K_0(\ofml K(H)) \cong \Z$; untuk itu, lihat~\cite{RordamLL:2000}, Akibat~6.4.2.)
\end{exam}

Berikutnya adalah sifat stabilitas penting dari funktor~$K_0$.
\begin{prop}\label{0062681} Jika $A$ aljabar-$C^*$, maka $K_0(A) \cong K_0(\M nA)$.
\end{prop}

%\begin{proof} Lihat~\cite{RordamLL:2000}, proposisi 4.3.8.   \ns \end{proof}





















\section{Limit Induktif}
Dalam bagian~\ref{sec_ind_limits}, kita mendefinisikan \emph{limit induktif} untuk sebarang kategori, dan dalam
bagian~\ref{sec_LFsps}, kita membatasi perhatian pada limit induktif ketat dari ruang vektor topologis
konveks lokal, yang berguna dalam konteks teori distribusi. Dalam bagian
ini, kita menelaah limit induktif dari barisan aljabar-$C^*$.

\begin{defn} Dalam sebarang kategori, suatu
 \index{induktif!barisan}%
 \index{barisan!induktif}%
\df{barisan induktif} adalah pasangan $(A,\phi)$, dengan $A = (A_j)$ barisan objek dan
$\phi = (\phi_j)$ barisan morfisme sedemikian sehingga $\phi_j\colon A_j \sto A_{j+1}$ untuk
setiap~$j$. Suatu
 \index{induktif!limit}%
 \index{limit!induktif}%
\df{limit induktif} (atau
 \index{langsung!limit}%
 \index{limit!langsung}%
\df{limit langsung}) dari barisan $(A,\phi)$ adalah pasangan $(L,\mu)$, dengan $L$ suatu objek dan
$\mu = (\mu_n)$ barisan morfisme $\mu_j\colon A_j \sto L$ yang memenuhi
 \begin{enumerate}
  \item[(i)] $\mu_j = \mu_{j+1} \circ \phi_j$ untuk setiap $j \in \N$, dan
  \item[(ii)] jika $(M,\lambda)$ adalah pasangan dengan $M$ suatu objek dan $\lambda = (\lambda_j)$
suatu barisan morfisme $\lambda_j\colon A_j \sto M$ yang memenuhi $\lambda_j = \lambda_{j+1}
\circ \phi_j$, maka terdapat morfisme tunggal $\psi\colon L \sto M$ sedemikian sehingga $\lambda_j
= \psi \circ \mu_j$ untuk setiap $j \in \N$.
 \end{enumerate}
Dengan sedikit penyalahgunaan bahasa, biasanya kita mengatakan bahwa $L$ adalah limit induktif dari barisan $(A_j)$ dan
 \index{limit@$\underrightarrow{\lim} A_j$ (limit induktif)}%
menulis $L = \underrightarrow{\lim} A_j$.

\vskip 15 pt

   \[\xymatrix@+15pt@C+25pt{&&&& L\ar@{-->}[dddd]^\psi \\
                   &&&&   \\
                  {\cdots}\ar[r] & A_j \ar[uurrr]^{\mu_j}\ar[ddrrr]^{\lambda_j}\ar[r]^(.6){\phi_j} &
                       A_{j+1}\ar[uurr]^{\mu_{j+1}}\ar[ddrr]^{\lambda_{j+1}}\ar[r] & {\cdots} &  \\
                   &&&&   \\
                   &&&& M
    }\]
\end{defn}

\begin{prop} Limit induktif dari barisan induktif (jika ada dalam suatu kategori) tunggal (hingga isomorfisme).
\end{prop}

\begin{prop} Setiap barisan induktif $(A,\phi)$ dari aljabar-$C^*$
 \index{cstaralg@aljabar-$C^*$!limit induktif}%
memiliki limit induktif. (Demikian pula setiap barisan induktif dari grup Abelian.)
\end{prop}

\begin{proof} Lihat \cite{Blackadar:2006}, II.8.2.1; \cite{RordamLL:2000}, proposisi 6.2.4; atau
\cite{Wegge-Olsen:1993}, lampiran~L.  \ns
\end{proof}

\begin{prop} Jika $(L,\mu)$ merupakan limit induktif dari barisan induktif $(A,\phi)$ dari
aljabar-$C^*$, maka $L = \clo{\bigcup {\mu_n}^\sto(A_n)}$ dan $\norm{\mu_m(a)} = \lim_{n \sto
\infty}\norm{\phi_{n,m}(a)} = \inf_{n \ge m}\norm{\phi_{n,m}(a)}$ untuk semua $m \in \N$ dan $a
\in A$.
\end{prop}

\begin{proof} Lihat \cite{RordamLL:2000}, proposisi 6.2.4.   \ns \end{proof}

\begin{defn} Suatu
 \index{aljabar-$C^*$ berdimensi hingga secara aproksimatif}%
 \index{aljabar-AF}%
 \index{hingga!berdimensi!secara aproksimatif}%
 \index{cstaralg@aljabar-$C^*$!aljabar-AF}%
\df{aljabar-$C^*$ berdimensi hingga secara aproksimatif} (singkatnya aljabar-AF) adalah limit induktif
dari suatu barisan aljabar-$C^*$ berdimensi hingga.
\end{defn}

\begin{exam} Jika $(A_n)$ merupakan barisan meningkat yang terdiri atas subaljabar-$C^*$ dari suatu aljabar-$C^*$ $D$
(dengan $\iota_n\colon A_n \sto A_{n+1}$ pemetaan inklusi untuk setiap~$n$), maka $(A,\iota)$
merupakan barisan induktif aljabar-$C^*$ yang limit induktifnya adalah $(B,j)$, dengan $B =
\clo{\bigcup_{n=1}^\infty A_n}$ dan $j_n\colon A_n \sto B$ pemetaan inklusi untuk setiap~$n$. % SOURCE-CORRECTION: CH17-C019
\end{exam}

\begin{exam} Barisan $\mathbf M_1 = \C \to^{\phi_1} \mathbf M_2 \to^{\phi_2} \mathbf M_3
\to^{\phi_3} \dots$ (dengan $\phi_n(a) = \diag(a,0)$ untuk setiap $n \in \N$ dan $a \in \mathbf
M_n$) merupakan barisan induktif aljabar-$C^*$ yang limit induktifnya adalah aljabar-$C^*$
$\ofml K(H)$ yang terdiri atas operator kompak pada ruang Hilbert terpisahkan berdimensi tak hingga~$H$. % SOURCE-CORRECTION: CH17-C020
\end{exam}

\begin{proof} Lihat \cite{RordamLL:2000}, bagian 6.4.  \ns \end{proof}

\begin{exam} Barisan $\Z \to^{\vc 1} \Z \to^{\vc 2} \Z \to^{\vc 3} \Z \to^{\vc 4} \dots$
(dengan $\vc n \colon \Z \sto \Z$ memenuhi $\vc n(1) = n$) merupakan barisan induktif grup
Abelian yang limit induktifnya adalah himpunan $\Q$ dari bilangan rasional.
\end{exam}

\begin{exam} Barisan $\Z \to^{\vc 2} \Z \to^{\vc 2} \Z \to^{\vc 2} \Z \to^{\vc 2} \dots$
merupakan barisan induktif grup Abelian yang limit induktifnya adalah himpunan bilangan rasional
diadik.
\end{exam}

Hasil berikut disebut
 \index{kekontinuan!dari $K_0$}%
 \index{K0@$K_0$!kekontinuan dari}%
\emph{sifat kekontinuan $K_0$}.

\begin{prop} Jika $(A,\phi)$ barisan induktif aljabar-$C^*$, maka
   \[ K_0\bigl(\,\underrightarrow{\lim}\, A_n\bigr) = \underrightarrow{\lim}\,K_0(A_n)\,. \]
\end{prop}

\begin{proof} Lihat \cite{RordamLL:2000}, teorema 6.3.2.   \ns \end{proof}

\begin{exam} Jika $H$ ruang Hilbert,
 \index{K0@$K_0(A)$!adalah $\Z$ ketika $A = \ofml K(H)$}%
maka $K_0(\ofml K(H)) \cong \Z$.
\end{exam}





















\section{Diagram Bratteli}
\begin{prop} Homomorfisme-$*\,$ tak nol dari $M_k$ ke $M_n$ ada hanya jika $n \ge k$; dalam
hal ini, semuanya tepat merupakan pemetaan berbentuk
   \[ \phi\colon a \mapsto u\,\diag(a, a, \dots, a, \vc 0)\, u^* \] % SOURCE-CORRECTION: CH17-C021
dengan $u$ matriks uniter. Di sini terdapat $m$ salinan $a$, dan $\vc 0$ adalah matriks nol berukuran $r \times
r$, dengan $n = mk + r$. Bilangan $m$ adalah
 \index{multiplisitas!dari homomorfisme-$*\,$}%
\df{multiplisitas} dari~$\phi$.
\end{prop}

\begin{proof} Lihat \cite{Power:1992}, akibat 1.3.  \ns \end{proof}

\begin{exam} Contoh homomorfisme-$*\,$ dari $M_2$ ke $M_7$ adalah
 \[  A =
  \begin{bmatrix} a & b \\ c & d  \end{bmatrix} \mapsto
    \begin{bmatrix}
      a & b & 0 & 0 & 0 & 0 & 0 \\
      c & d & 0 & 0 & 0 & 0 & 0 \\
      0 & 0 & a & b & 0 & 0 & 0 \\
      0 & 0 & c & d & 0 & 0 & 0 \\
      0 & 0 & 0 & 0 & 0 & 0 & 0 \\
      0 & 0 & 0 & 0 & 0 & 0 & 0 \\
      0 & 0 & 0 & 0 & 0 & 0 & 0
    \end{bmatrix} = \diag(A,A, \vc 0)\]
 dengan $m = 2$ dan $r = 3$.
\end{exam}

\enlargethispage{\baselineskip}

\begin{prop}\label{0068117} Setiap aljabar-$C^*$ berdimensi hingga $A$ isomorfik dengan
jumlah langsung aljabar matriks
   \[ A \simeq M_{k_1} \oplus \dots \oplus M_{k_r}.\]
Misalkan $A$ dan $B$ aljabar-$C^*$ berdimensi hingga, sehingga
   \[ A \simeq M_{k_1} \oplus \dots \oplus M_{k_r} \qquad \text{dan}
                          \qquad B \simeq M_{n_1} \oplus \dots \oplus M_{n_s} \]
dan misalkan $\phi\colon A \sto B$ homomorfisme-$*\,$ beridentitas. Maka $\phi$
ditentukan (hingga ekuivalensi uniter dalam~$B$) oleh matriks $s \times r$ $\vc m = \bigl[
m_{ij}\bigr]$ yang berentri bilangan bulat tak negatif sedemikian sehingga % SOURCE-CORRECTION: CH17-C022
   \begin{equation}\label{0068117i}
    \vc m\, \vc k = \vc n.
   \end{equation}
\emph{Di sini $\vc k = (k_1, \dots,k_r)$, $\vc n = (n_1, \dots, n_s)$, dan bilangan $m_{ij}$ adalah
multiplisitas pemetaan}
 \[\xymatrix@1{M_{k_j} \ar[r] & A \ar[r]^\phi & B \ar[r] & M_{n_i}
 }\]
\end{prop}

\begin{proof} Lihat \cite{Davidson:1996}, teorema III.1.1.  \ns \end{proof}

\begin{defn} Misalkan $\phi$ seperti dalam proposisi~\ref{0068117} sebelumnya. Suatu
 \index{diagram Bratteli}%
 \index{diagram!Bratteli}%
\df{diagram Bratteli} untuk $\phi$ terdiri atas dua baris (atau kolom) simpul yang diberi label
$k_j$ dan $n_i$, bersama $m_{ij}$ sisi yang menghubungkan $k_j$ ke $n_i$ (untuk $1 \le j \le
r$ dan $1 \le i \le s$).
\end{defn}

\begin{exam} Misalkan $\phi\colon \C \oplus \C \sto M_4 \oplus M_3$ diberikan oleh
 \[ (\lambda, \mu) \mapsto \left(
      \begin{bmatrix}
        \lambda & 0 & 0 & 0 \\
        0 & \lambda & 0 & 0 \\
        0 & 0 & \lambda & 0 \\
        0 & 0 & 0 & \mu
      \end{bmatrix}\, , \,
      \begin{bmatrix}
        \lambda & 0 & 0 \\
        0 & \lambda & 0 \\
        0 & 0 & \mu
      \end{bmatrix}\right)\, .\]
Maka $m_{11} = 3$, $m_{12} = 1$, $m_{21} = 2$, dan $m_{22} = 1$. Perhatikan bahwa
   \[\vc m \,\vc k = \begin{bmatrix} 3 & 1 \\ 2 & 1 \end{bmatrix}\,
             \begin{bmatrix} 1 \\ 1 \end{bmatrix} =
             \begin{bmatrix} 4 \\ 3 \end{bmatrix} = \vc n.\]
Suatu diagram Bratteli untuk $\phi$ adalah
    \[\xymatrix@=90pt{1 \ar@<.7ex>[r]\ar[r]\ar@<-.7ex>[r]\ar@<.5ex>[dr]\ar@<-.5ex>[dr] & 4 \\
                1 \ar[ur]\ar[r] & 3
     }\]
\end{exam}

\begin{exam} Misalkan $\phi\colon \C \oplus M_2 \sto M_3 \oplus M_5 \oplus M_2$
diberikan oleh
 \[ (\lambda, b) \mapsto \left(
    \begin{bmatrix}
        \lambda & 0  \\
            0 &  b
      \end{bmatrix}\, , \,
      \begin{bmatrix}
        \lambda & 0 & 0  \\
              0 & b & 0  \\
              0 & 0 &  b
      \end{bmatrix} \,,\, b \, \right).
    \]
Maka $m_{11} = 1$, $m_{12} = 1$, $m_{21} = 1$, $m_{22} = 2$, $m_{31} = 0$, dan $m_{32} = 1$.
Perhatikan bahwa
 \[\vc m \,\vc k = \begin{bmatrix} 1 & 1 \\ 1 & 2 \\ 0 & 1 \end{bmatrix}\,
             \begin{bmatrix} 1 \\ 2 \end{bmatrix} =
             \begin{bmatrix} 3 \\ 5 \\2 \end{bmatrix} = \vc n.\]
Suatu diagram Bratteli untuk $\phi$ adalah
 \[\xymatrix@=20pt@C+30pt{&& 3                    \\
                   1 \ar[urr]\ar[drr] &&   \\
                   && 5                    \\
         2\ar[uuurr]\ar@<+.5ex>[urr]\ar@<-.5ex>[urr]\ar[drr] && \\
                   && 2
   }\]
\end{exam}

\begin{exam} Misalkan $\phi\colon M_3 \oplus M_2 \oplus M_2 \sto M_9 \oplus M_7$
diberikan oleh

\vskip 10 pt

 \[ (a,b,c) \mapsto \left(
    \begin{bmatrix}
            a & \vc 0 & \vc 0 \\
        \vc 0 &     a & \vc 0 \\
        \vc 0 & \vc 0 & \vc 0
      \end{bmatrix}\, , \,
      \begin{bmatrix}
            a & \vc 0 & \vc 0  \\
        \vc 0 &     b & \vc 0  \\
        \vc 0 & \vc 0 &     c
      \end{bmatrix}\, \right)\, .\]
Maka $m_{11} = 2$, $m_{12} = 0$, $m_{13} = 0$, $m_{21} = 1$, $m_{22} = 1$, dan $m_{23} = 1$.
Perhatikan bahwa kali ini terjadi sesuatu yang ganjil: kita mempunyai $\vc m\,\vc k \ne \vc n$:
 \[\vc m \,\vc k = \begin{bmatrix} 2 & 0 & 0 \\ 1 & 1 & 1 \end{bmatrix}\,
             \begin{bmatrix} 3 \\ 2 \\ 2 \end{bmatrix} =
             \begin{bmatrix} 6 \\ 7 \end{bmatrix} \le
             \begin{bmatrix} 9 \\ 7 \end{bmatrix} = \vc n.\]
Apa masalahnya di sini? Nah, hasil yang dinyatakan sebelumnya berlaku untuk homomorfisme-$*\,$
\emph{beridentitas}, sedangkan $\phi$ ini \emph{tidak} beridentitas. Secara umum, inilah hasil terbaik
yang dapat kita harapkan: $\vc m\,\vc k \le \vc n$.

\noindent Berikut diagram Bratteli yang dihasilkan untuk $\phi$:
   \[\xymatrix@=+20pt@C+20pt@R-7pt{
        3\ar@<+.5ex>[drr]\ar@<-.5ex>[drr]\ar[dddrr] && \\
                            && 9                   \\
                           2\ar[drr] &&            \\
                            && 7                   \\
                           2\ar[urr] &&
   }\]
\end{exam}

\begin{exam} Jika $K$ himpunan Cantor, maka $\fml C(K)$ merupakan aljabar-AF.
 \index{AF@aljabar-AF!$\fml C(K)$ merupakan aljabar-AF ($K$ himpunan Cantor)}%
Untuk melihatnya, tuliskan $K$ sebagai irisan dari suatu keluarga menurun yang terdiri atas himpunan bagian tertutup $K_j$, yang masing-masing
terdiri atas $2^j$ subinterval tertutup yang saling lepas dari $[0,1]$. Untuk setiap $j \ge 0$, misalkan
$A_j$ subaljabar yang terdiri atas fungsi-fungsi dalam $\fml C(K)$ yang konstan pada setiap
interval penyusun~$K_j$. Jadi $A_j \simeq \C^{2^j}$ untuk setiap~$j \ge 0$. Pembenaman
 \[\phi_j \colon A_j \sto A_{j+1}\colon
      (a_1, a_2, \dots, a_{2^j}) \mapsto
        (a_1, a_1, a_2, a_2,\dots, a_{2^j}, a_{2^j})\]
memecah setiap proyeksi minimal menjadi jumlah dua proyeksi minimal. Untuk $\phi_0$, matriks
$\vc m_0$ dari ``multiplisitas parsial'' yang bersesuaian adalah $\begin{bmatrix} 1 \\ 1
\end{bmatrix}$; matriks $\vc m_1$ yang bersesuaian dengan $\phi_1$ adalah
$\begin{bmatrix} 1 & 0\\ 1 & 0 \\ 0 & 1 \\ 0 & 1 \end{bmatrix}$; dan seterusnya. Jadi diagram Bratteli
untuk limit induktif $\fml C(K) = \underrightarrow{\lim}A_j$ adalah:

\vskip 20 pt

 \[\xymatrix@=+7pt@C+20pt@R-7pt{
     &&& 1 & \\
     && 1\ar[ur]\ar[dr] && \\
     &&& 1 & \\
     & 1\ar[uur]\ar[ddr] &&& \\
     &&& 1 & \\
     && 1\ar[ur]\ar[dr] && \\
     &&& 1 & \\
     1\ar[uuuur]\ar[ddddr] &&&& \cdots \\
     &&& 1 & \\
     && 1\ar[ur]\ar[dr] && \\
     &&& 1 & \\
     & 1\ar[uur]\ar[ddr] &&& \\
     &&&1 & \\
     && 1\ar[ur]\ar[dr] && \\
     &&& 1 &
   }\]
\end{exam}

\begin{exam} Berikut suatu contoh dari apa yang disebut
 \index{aljabar CAR}%
 \index{cstaralg@aljabar-$C^*$!aljabar CAR}%
aljabar CAR (CAR = Relasi Antikomutasi Kanonik). Untuk $j \ge 0$, misalkan $A_j = % SOURCE-CORRECTION: CH17-C023
M_{\spsp{2}j}$ dan
 \[\phi_j\colon A_j \to A_{j+1}\colon \vc a \mapsto
      \begin{bmatrix} \vc a & 0 \\ 0 & \vc a \end{bmatrix}.\]
``Matriks'' multiplisitas $\vc m$ untuk setiap $\phi_j$ hanyalah matriks $1 \times 1$~$[2]$.
Kita melihat bahwa untuk setiap $j$
 \[\vc m\, \vc k(j) = [2]\,[2^j] = [2^{j+1}] = \vc n(j).\]
Ini menghasilkan diagram Bratteli berikut

\vskip 15pt

 \[\xymatrix{1 \ar@<.5ex>[r]\ar@<-.5ex>[r] & 2 \ar@<.5ex>[r]\ar@<-.5ex>[r] & 4
        \ar@<.5ex>[r]\ar@<-.5ex>[r] & 8 \ar@<.5ex>[r]\ar@<-.5ex>[r] & 16
        \ar@<.5ex>[r]\ar@<-.5ex>[r] & 32
        \ar@<.5ex>[r]\ar@<-.5ex>[r]& {\cdots}
   }\]
\vskip 15pt
\noindent untuk limit induktif $C = \underrightarrow{\lim}A_j$.

\textbf{Namun}, jangkauan $\phi_j$ termuat dalam subaljabar $B_{j+1} \simeq M_{2^j}
\oplus M_{2^j}$ dari~$A_{j+1}$. (Ambil $B_0 = A_0 = \C$.) Jadi, untuk $j \in \N$, kita dapat memandang
$\phi_j$ sebagai pemetaan dari $B_j$ ke~$B_{j+1}$:
 \[ \phi_j\colon(\vc b, \vc c) \mapsto
    \left(\begin{bmatrix} \vc b & 0 \\ 0 & \vc c \end{bmatrix},
          \begin{bmatrix} \vc b & 0 \\ 0 & \vc c \end{bmatrix}
              \right)\,.\]
Sekarang matriks multiplisitas $\vc m$ untuk setiap $\phi_j$ adalah $\begin{bmatrix} 1 & 1 \\ 1 & 1
\end{bmatrix}$, dan kita melihat bahwa untuk setiap $j$
 \[\vc m\, \vc k(j) =
   \begin{bmatrix} 1 & 1 \\ 1 & 1 \end{bmatrix}\,
        \begin{bmatrix} 2^{j-1} \\ 2^{j-1} \end{bmatrix} =
   \begin{bmatrix} 2^j \\ 2^j \end{bmatrix} = \vc n(j).\]
Karena $C = \underrightarrow{\lim}A_j = \underrightarrow{\lim}B_j$, kita memperoleh diagram Bratteli kedua
(yang cukup berbeda) untuk aljabar-AF~$C$.
   \[\xymatrix@-10pt{&& 1 \ar[rr]\ar[rrdd] && 2 \ar[rr]\ar[rrdd] && 4 \ar[rr]\ar[rrdd]
                                && 8 \ar[rr]\ar[rrdd] && 16 \ar[rr]\ar[rrdd]&& \cdots\\
              1 \ar[urr]\ar[drr]&&&&&&&&&&  \\
             && 1 \ar[rr]\ar[rruu] && 2 \ar[rr]\ar[rruu] && 4 \ar[rr]\ar[rruu]
                                && 8 \ar[rr]\ar[rruu] && 16 \ar[rr]\ar[rruu]&& \cdots \\
   }\]
\end{exam}

\begin{exam} Ini adalah
 \index{aljabar Fibonacci}%
 \index{cstaralg@aljabar-$C^*$!aljabar Fibonacci}%
\df{aljabar Fibonacci}. Untuk $j \in \N$, definisikan barisan $\bigl(p_j\bigr)$ dan
$\bigl(q_j\bigr)$ dengan relasi rekursi yang dikenal:
 \begin{align*}
       p_1 = &q_1 = 1, \\
       p_{j+1} = &p_j + q_j, \text{ dan } \\
       q_{j+1} &= p_j.
 \end{align*}
Untuk semua $j \in \N$, misalkan $A_j = M_{p_j} \oplus M_{q_j}$ dan

\vskip 10 pt

 \[\phi_j\colon A_j \sto A_{j+1}\colon (a,b) \mapsto
     \left(\begin{bmatrix} a & 0 \\ 0 & b
          \end{bmatrix},  a \right). \]
Matriks multiplisitas $\vc m$ untuk setiap $\phi_j$ adalah $\begin{bmatrix} 1 & 1 \\ 1 & 0
\end{bmatrix}$, dan untuk setiap $j$
 \[\vc m\, \vc k(j) =
   \begin{bmatrix} 1 & 1 \\ 1 & 0 \end{bmatrix}\,
        \begin{bmatrix} p_j \\ q_j \end{bmatrix} =
   \begin{bmatrix} p_j + q_j \\ p_j \end{bmatrix} =
        \begin{bmatrix} p_{j+1} \\ q_{j+1} \end{bmatrix} =
                          \vc n(j).\]
Diagram Bratteli yang dihasilkan untuk $F = \underrightarrow{\lim}A_j$ adalah
   \[\xymatrix@-10pt{&& 1 \ar[rr]\ar[rrdd] && 2 \ar[rr]\ar[rrdd] &&
         3 \ar[rr]\ar[rrdd] && 5 \ar[rr]\ar[rrdd] && 8 \ar[rr]\ar[rrdd]&& \cdots\\
              1 \ar[urr]\ar[drr]&&&&&&&&&&  \\
             && 1 \ar[rruu] && 1 \ar[rruu] && 2 \ar[rruu] &&
              3 \ar[rruu] && 5 \ar[rruu]&& \cdots \\
   }\]
\end{exam}













\endinput
